\section{Discussion}
\label{sec:discussion}

\subsection{Why CBAM and BiomedCLIP are Complementary}
% TODO: prose version of the table in MOTIVATION_AND_DESIGN.md §5.
The 3D CBAM backbone and the per-slice BiomedCLIP encoder operate on
informationally orthogonal axes. CBAM provides volumetric spatial
localisation, learned end-to-end on RSNA's stenosis task. BiomedCLIP
provides language-grounded semantic features, learned by contrastive
pretraining on 15M PubMed image–text pairs. The projection MLP fuses
both into a single 512-d embedding aligned with the CLIP text space.

\subsection{Selective Transfer in Phase 3}
% Honest framing of the win/loss split — turn it into a positive
% characterisation result.
The Phase~3 zero-shot pattern is consistent with the complementarity
story: the Hybrid wins on disc-related labels (Disc\_narrowing,
Disc\_bulging, Disc\_herniation) where volumetric geometry across
slices carries discriminative signal; ``Naked BiomedCLIP'' wins on
vertebra/endplate labels (Modic, UP/LOW endplate, Spondylolisthesis)
that are visually distinctive in 2D and well-represented in PubMed.
This is not a failure mode but a useful characterisation: the hybrid
adds value when 3D context matters and the projection MLP's RSNA
specialisation is aligned with the target task. We report this
transparently rather than smoothing the mean.

\subsection{Why Foundation-Model-Agnostic}
% Plug-in story: BiomedCLIP can be swapped for MedSigLIP/RadFM/etc.
Because the integration uses a frozen contrastive encoder as a black
box, BiomedCLIP can be replaced by any future medical contrastive
encoder (e.g.\ MedSigLIP from~\cite{sebro2025medgemma},
RAD-DINO~\cite{perezgarcia2024raddino}) with only a vocabulary remap.
The 3D CBAM specialist branch likewise admits stronger alternatives
without changing the cosine-similarity head.

\subsection{Limitations}
\begin{itemize}
  \item Single-seed reporting; multi-seed mean$\pm$std is planned for
        the camera-ready.
  \item Validation is patient-level held-out within each dataset;
        SPIDER zero-shot is the strongest cross-cohort test we have
        but is still single-site.
  \item The projection MLP induces a label-space bias toward the
        training task (here, RSNA stenosis). Multi-domain projection
        heads are an obvious next step.
\end{itemize}
